\documentclass[10pt]{article}

% graphicx package, useful for including eps, pdf, and raster graphics
\usepackage{graphicx}
\DeclareGraphicsExtensions{.png,.png,.jpg,.pdf}

% basic packages
\usepackage{fancyhdr}
\usepackage{color}
\usepackage{parskip}
\usepackage{float}
\usepackage{microtype}
\usepackage{url}
\usepackage{hyperref}
\usepackage{booktabs}
\usepackage{makecell}
\usepackage{multirow}
\usepackage{pbox}
\usepackage{enumitem}

% fonts
\usepackage{fontspec}
\setmainfont{DroidSans}[
  Extension   = .ttf,
  UprightFont = *,
  BoldFont    = *-Bold ]
\setmonofont{DroidSansMono}[
  Extension   = .ttf,
  UprightFont = * ]

\hypersetup{colorlinks,urlcolor=blue}
\urlstyle{same}

% page layout
\usepackage[top=1in, bottom=1in, left=0.75in, right=0.75in, includefoot, heightrounded]{geometry}

% Useful for seeing the layout as you tweak it.
%\usepackage{showframe}

\setlength{\headheight}{1.1in}
\pagestyle{fancy}
\fancyhf{}

\lhead{\includegraphics[height=1in,clip]{logo}}
\rhead{\pbox[b][1in][c]{\textwidth}{
    \small
    \textbf{Test Performed By}\\
    Northwest Genomics Center\\
    University of Washington School of Medicine\\
    Genome Sciences Box 355065\\
    3720 15th Ave NE\\
    Seattle, WA 98195-5065}}

\newcommand{\PageLine}{\rule{\textwidth}{0.25mm}}

\newcommand{\link}[2]{\href{#1}{#2}\footnote{#1}}

\begin{document}

\begin{center}
\Large
\textbf{Warbixinta Shaybaadhka NWGC SARS-CoV-2 (COVID-19) Qual Jawaab Ce­linta Polymerase (Polymerase chain reaction, PCR)}
\end{center}

\bigskip

\begin{description}[font=\normalfont,align=left,labelwidth=16em]
\item [Daryeelka Kaqayb Galaha] \textbf{\VAR{pat_name|e}}
\item [Participant Date of Birth] \textbf{\VAR{birth_date|e}}
\item [Aqoonsiga Muunada] \textbf{\VAR{qrcode|e}}
\item [Taariikhda Muunada La Gudbiyay] \textbf{\VAR{collect_ts|e}}
\item [Taariikhda natiijada la ansixiyay]
  %- if status_code == "never-tested"
  \textbf{Never tested}
  %- else
  \textbf{\VAR{result_ts|e}}
  %- endif
\item [Nooca Muunada] Duufka, ka qaadid
\end{description}

\PageLine

Natiijada NWGC SARS-CoV-2 (COVID-19) Qual PCR

%- if status_code == "negative"
\textbf{Lama ogaanin}\\
%- elif status_code == "positive"
\textbf{SARS-CoV-2 La Ogaaday (laga heley)}\\
%- elif status_code == "inconclusive"
\textbf{Aan Joogto ahayn}\\
%- elif status_code == "never-tested"
\textbf{La baadhi kari waayay}\\
%- else
\VAR{0/0} %# poor man's raise
%- endif

\PageLine

Fadlan Ogsoonaw:

\begin{itemize}
\item

  \textbf{Muunada noocan ah iyo shaybaadhka loo isticmaalayo mashruucan waxaa
  loo sameeyay oo uu qayb ka ahaa barnaamijka kormeerka caafimaadka oo la iska
  kaashanayo Degmada King iyo hay'ado kale.} Shaybaadhka aad ka heshay ma
  badalayo daryeelka caafimaadka shakhsi ahaaneed ee daryeelka caafimaad
  bixiyaha.

\item

  \textbf{Silsilada Muunada Haynta:} muunada la tijaabiyay qofka ayaa iska
  qaaday, calaamadiyay, oo markaa usoo gudbiyay Dalada Qiimaynta Coronavirus
  Seattle (Seattle Coronavirus Assessment Network, SCAN) boosta ama adeega
  keenista. Lama qaadin, calaamadin, ama muu qaadi xirfadle daryeelka caafimaad
  ah. Tani waxa ay saamayn kartaa natiijada shaybaadhka iyo natiijo been ah
  markaa lagama reebi karo. La socodka shaybaadhka xarunta caafimaadka lama
  damaanad qaadi karo.

\item

  \textbf{Aalada qaadida muunada:} hadii muunada ayna lahayn Maamulka
  Cuntooyinka iyo Dawooyinka/Xarumaha Kahortaga iyo Xakamaynta Xanuunada (Food
  and Drug Administration, FDA/ Centers for Disease Control and Prevention, CDC)
  Ogolaanshaha Isticmaalka Xaalada Degdeg (Emergency Use Authorization, EUA)
  nooca muunada ee la ansixiyay. Markaa, lama ogaanin/natiijo aan laga helin
  waxaa dhici karta inayna reebeynin caabuqa SARS-CoV-2.

\item

  \textbf{Natiijada ``Laga heley'' waxaa dhici karta inay cadayso Washington
  State Department of Health}. Natiijadan waa in laga wada hadlaa daryeelka
  caafimaad bixiye hadii aad qabto su'aal.

\item{
  \textbf{Hadii shaybaadhkaagu uu yidhaahdo ``Laguma arag'' waxa ay ka dhigan
  tahay in SARS-CoV-2 ayna ku arag shaybaadhkeeda.}

Ma jiro shaybaadh wanaagsan oo waxaa jirta fursad yar oo ah inaad qabto fayruska
SARS-CoV-2 kaas oo keena COVID-19 iyada oo shaybaadhku sheegayo inaan lagaa
helin.  Shaybaadhkani ma isticmaalo qaabka caadiga ah ee uu ku qaado muunada
daryeelka caafimaad bixiyuhu, oo muunada waxaa lagu dirayaa boosta ama waxaa loo
gaynayaa adeega, laguma qaadin oo manay soo qaaday daryeelka caafimaad bixiye.
Waxa uu doono ha noqdo natiijada shaybaadhku, waxaa wanaagsan inaad joogto
guriga hadii aad xanuunsato, iyo inaad iska ilaaliso inaad dadka kale isasoo
gaadhaan. Hadii aad qabto mushkilad kale oo caafimaad ama aad xanuusanayso ama
aad ka warwarsan tahay inaad xanuunsato, waa inaad la xidhiidhaa daryeelka
caafimaad bixiye.  Hadii raadka xanuunkaagu uu joogto noqdo, waa inaad kala
hadashaa sidii aad u mari lahayn shaybaadh kale oo COVID-19 ah daryeelka
caafimaad bixiyahaaga.
}

\item

  \textbf{Hadii shaybaadhkaaga ay kasoo baxdo natiijo ``Aan dhamaystirnayn''}
  waxa ay ka dhigan tahay in natiijada shaybaadhkaaga SARS-CoV-2 aan la hubin.

  Ma jiro wax shaybaadh ah oo boqolkiiba boqol ah oo waxa kasoo bixi karta
  natiijo aan dhamaystirnayn taas oo ka dhigan inaad qabto fayruska SARS-CoV-2
  ee keena COVID-19.  Shaybaadhkan looma adeegsado qaabka caadiga ah ee uu ku
  qaado muunada daryeelka caafimaad bixiye, oo muunadan waxaa lagu soo diray
  boosta ama waxaa loo keeney adeega, e muu qaadin ama muusoo qaadin daryeelka
  caafimaad bixiye.  Waxa uu doono ha noqdo natiijada shaybaadhku, waxaa
  wanaagsan inaad joogto guriga hadii aad xanuunsato, iyo inaad iska ilaaliso
  inaad dadka kale isasoo gaadhaan. Hadii aad qabto mushkilad kale oo caafimaad
  ama aad xanuusanayso ama aad ka warwarsan tahay inaad xanuunsato, waa inaad la
  xidhiidhaa daryeelka caafimaad bixiye.  Hadii raadka xanuunkaagu uu joogto
  noqdo, waa inaad kala hadashaa sidii aad u mari lahayn shaybaadh kale oo
  COVID-19 ah daryeelka caafimaad bixiyahaaga.

\item

  \textbf{Hadii shaybaadhkaagu uu ku qoran yahay ``Waa la baadhi kari waayay''}
  waxa ay ka dhigan tahay inayna jirin muunada ku filan ama inuu khalad ka
  dhacay shaybaadhka.

  Maadaama oo la samayn kari waayay shaybaadhka, waxaa dhici karta inaad qabto
  SARS-CoV-2 fayruska ah ee keena COVID-19.  Shaybaadhkani ma isticmaalo qaabka
  caadiga ah ee uu ku qaado muunada daryeelka caafimaad bixiyuhu, oo muunada
  waxaa lagu dirayaa boosta ama waxaa loo gaynayaa adeega, laguma qaadin oo
  manay soo qaaday daryeelka caafimaad bixiye. Waxa uu doono ha noqdo natiijada
  shaybaadhku, waxaa wanaagsan inaad joogto guriga hadii aad xanuunsato, iyo
  inaad iska ilaaliso inaad dadka kale isasoo gaadhaan. Hadii aad qabto
  mushkilad kale oo caafimaad ama aad xanuusanayso ama aad ka warwarsan tahay
  inaad xanuunsato, waa inaad la xidhiidhaa daryeelka caafimaad bixiye.  Hadii
  raadka xanuunkaagu uu joogto noqdo, waa inaad kala hadashaa sidii aad u mari
  lahayn shaybaadh kale oo COVID-19 ah daryeelka caafimaad bixiyahaaga.

\item

  \textbf{Shaybaadh laga heley (waxa uu muujinayaa in ``LAGU ARKAY'' ama
  ``Fayruska La Arkay'') waxaana ay gaadhsiin doontaa Waaxda Caafimaadka gobolka
  iyo degmada ee munaasabka ah xarunta baadhaysaayi.} LAGAAGAMA baahna inaad
  aftaadu aad gaadhsiiso warbixinta shaybaadhka. Caafimaadka Shacabka waxa aad
  ka wada socon kartaan hadii lagaa helo novel coronavirus si ay markaa kaaga
  waraystaan xanuunkaaga iyo dadka aad isku dhawateen.

\item

  \textbf{Waxa aad daabacan kartaa ama iimeel ku diran kartaa natiijada
  shaybaadhka daryeelka caafimaad bixiyaha.}

\end{itemize}

\bigskip
\textbf{Qaabka NWGC SARS-CoV-2 (COVID-19) Qual PCR:}

Qaabkan waxaa loogu talo geley in lagu ogaado SARS-CoV-2 iyada oo la adeegsanayo
qoraalka PCR. Qaabkan waxa uu ku tiirsan yahay waynaysada qaybta jeermiska Orf1B
iyo S ee gudaha bilaa aadanka iyada oo la adeegsanayo waynaysada RNaseP.
Shaybaadhka waxaa samaynaya oo qaabka loo samaynayaa waxa sheegeysa University
of Washington’s Northwest Genomics Center.  Shaybaadhka uu sameeyay shaybaadhku
waa la ansixiyay oo waxa loo gudbiyay Maamulka Cuntada iyo Dawada Maraykanka si
loogu sameeyo Ogolaanshe Isticmaalka xaalada Degdega ah. FDA ayaa si madaxbanaan
u naqdin doonta oo ansixin doonta shaybaadhka taas oo markaa la sugayo.
Shaybaadhkan waxaa lagu xaqiijiyay oo uu hoos yimaadaa Waxka Badalka Hormarka
Shaybaadhka Xarunta Caafimaad (Clinical Laboratory Improvement Amendments, CLIA)
taas oo loogu qalmo in lagu sameeyo shaybaadh murugsan oo xarunta caafimaadka
ah ..

Wixii macluumaad dheeraad ah, fadlan waxa aad eegtaa linkiga ay kusiiyeen Waaxda
Gobolka Washington ee Caafimaadku:

\begin{itemize}
\item

  \link{https://www.doh.wa.gov/Portals/1/Documents/1600/coronavirus/COVIDcasepositive.pdf}{Bukaanka
  looga shakiyo ama lagu xaqiijiyo inuu qabo COVID-19}

\item

  \link{https://www.doh.wa.gov/Portals/1/Documents/1600/coronavirus/COVIDexposed.pdf}{Bukaanka
  ay isasoo gaadheen dad qaba oo lagu xaqiijiyay COVID-19}

\item

  \link{https://www.doh.wa.gov/Portals/1/Documents/1600/coronavirus/COVIDconcerned.pdf}{Bukaan
  ayna isasoo gaadhin cid qabta oo leh calaamado COVID-19}

\end{itemize}

Khayraad dheeraad ah ayaa laga heleyaa websaytyada Waaxda Caafimaadka Gobolka
Washington iyo CDC:

\begin{itemize}
\item

  \url{https://www.doh.wa.gov}

\item

  \url{https://www.cdc.gov/coronavirus/2019-ncov/index.html}

\end{itemize}

\end{document}
